\section*{Bab \ref{ch:b3:galois} --- Perluasan Lapangan dan
Teori Galois}

\begin{solution}{exo:b3:galois:1}
$\sqrt3 \notin \Q(\sqrt2)$: sebab dari $\sqrt3 = a + b\sqrt2$ (dengan $a, b
\in \Q$), pengkuadratannya memberikan $3 = a^2 + 2b^2 + 2ab\sqrt2$, sehingga $ab =
0$; lalu $b = 0$ membuat $\sqrt3$ rasional, sedangkan $a = 0$ memberikan $\sqrt6 = 2b
\in \Q$ --- keduanya salah (lewat argumen faktorisasi prima yang baku).
Karena itu $[\Q(\sqrt2,\sqrt3) : \Q(\sqrt2)] = 2$ dan \omterm{thm:b3:galois:tower}{hukum menara}
memberikan \omterm{def:b3:galois:extension}{derajat} $4$.

Misalkan $\gamma = \sqrt2 + \sqrt3$. Maka $\gamma^2 = 5 + 2\sqrt6$ dan
$(\gamma^2 - 5)^2 = 24$: jadi $\gamma$ menghapus $X^4 - 10X^2 + 1$.
Lebih lanjut $\gamma^3 = 11\sqrt2 + 9\sqrt3$, sehingga $\gamma^3 - 9\gamma
= 2\sqrt2$: berarti $\sqrt2 \in \Q(\gamma)$, lalu $\sqrt3 = \gamma -
\sqrt2 \in \Q(\gamma)$: jadi $\Q(\gamma) = \Q(\sqrt2,\sqrt3)$, yang
berderajat $4$. Polinomial kuartik penghapus itu, karena berderajat sama dengan
\omterm{def:b3:galois:algebraic}{polinomial minimalnya}, \emph{adalah} \omterm{def:b3:galois:algebraic}{polinomial minimalnya}: $X^4 -
10X^2 + 1$ (khususnya ia \omterm{def:b3:rings:divisibility}{tak tereduksi} atas $\Q$).
\end{solution}

\begin{solution}{exo:b3:galois:2}
Ketiga akar $X^3 - 2$ di $\C$ adalah $\alpha, j\alpha,
j^2\alpha$ dengan $j = \eu^{2\iu\pi/3}$. Pemetaan $\alpha \mapsto
j\alpha$ mendefinisikan pembenaman-$\Q$ $\Q(\alpha) \to \C$
(\cref{thm:b3:galois:simple}: keduanya membangun perluasan berderajat $3$
dengan polinomial minimal yang sama), yang petanya $\Q(j\alpha)
\not\subseteq \R$ berbeda dari $\Q(\alpha) \subseteq \R$: jadi
$\Q(\alpha)/\Q$ tidak normal. \omterm{thm:b3:galois:splitting}{Lapangan pemecahnya} adalah $L =
\Q(\alpha, j)$, dengan $[L:\Q] = [L:\Q(\alpha)]\,[\Q(\alpha):\Q] =
2 \cdot 3 = 6$ (sebab $j$ memenuhi $X^2 + X + 1$, yang \omterm{def:b3:rings:divisibility}{tak tereduksi} atas
lapangan real $\Q(\alpha)$). Sebuah automorfisma $\Q(\alpha)$
harus mengirim $\alpha$ ke akar $X^3 - 2$ yang berada \emph{di dalam}
$\Q(\alpha) \subseteq \R$: hanya $\alpha$ yang memenuhinya, sehingga
$\operatorname{Aut}_\Q(\Q(\alpha)) = \{\mathrm{id}\}$, yang berorde
$1 < 3$.
\end{solution}

\begin{solution}{exo:b3:galois:3}
Polinomial $X^2 + 1$ tak berakar di $\mathbb F_3$ (sebab $0, 1, 2 \mapsto 1, 2,
2$), sehingga $\mathbb F_9 = \mathbb F_3[X]/(X^2+1)$ merupakan lapangan dengan
$9$ unsur; tulis $\omega = \bar X$, sehingga $\omega^2 = -1$. Grup
$\mathbb F_9^\times$ siklik berorde $8$; unsur $\omega$ berorde
$4$, tetapi $1 + \omega$ berhasil: $(1+\omega)^2 = 1 + 2\omega +
\omega^2 = 2\omega$, dan $(1+\omega)^4 = 4\omega^2 = \omega^2 = -1
\ne 1$: jadi berorde $8$.

Atas $\mathbb F_2$ --- berderajat $1$: $X$, $X + 1$; berderajat $2$:
$X^2 + X + 1$ (sebab ketiga polinomial kuadrat lainnya berakar); berderajat
$3$: $X^3 + X + 1$ dan $X^3 + X^2 + 1$ (tak berakar di $\mathbb
F_2$; keenam polinomial kubik lainnya berakar). Pemeriksaannya:
\[
(X^3{+}X{+}1)(X^3{+}X^2{+}1) = X^6 + X^5 + X^4 + X^3 + X^2 + X
+ 1,
\]
dan $X(X{+}1)(X^6 + \dots + 1) = X(X^7 + 1) = X^8 + X = X^8 -
X$ atas $\mathbb F_2$ --- tepat berisi \omterm{def:b3:rings:divisibility}{unsur tak tereduksi} yang \omterm{def:b3:galois:extension}{derajatnya}
membagi $3$, seperti diramalkan \cref{exo:b3:galois:6} (\omterm{def:b3:galois:extension}{derajat} $2$
absen sebab $2 \nmid 3$).
\end{solution}

\begin{solution}{exo:b3:galois:4}
(a) Modulo $7$: pangkat $3$ adalah $3, 2, 6, 4, 5, 1$: berorde
$6$, jadi sebuah pembangun; akar primitifnya adalah $3^k$ dengan
$\gcd(k, 6) = 1$: yakni $3$ dan $3^5 = 5$. Modulo $11$: pangkat $2$:
$2, 4, 8, 5, 10, 9, 7, 3, 6, 1$: jadi sebuah pembangun; akar primitifnya
$2^k$ dengan $\gcd(k, 10) = 1$: yakni $2, 2^3 = 8, 2^7 = 7, 2^9 = 6$.

(b) Tulis $x = g^k$ dengan $g$ sebuah pembangun
(\cref{thm:b3:galois:cyclic}). Maka $x$ merupakan kuadrat jika dan hanya jika $k$
genap (sebab kuadratnya adalah $g^{2l}$, dan $g^{2l} = g^{k}$ jika dan hanya jika $k
\equiv 2l \bmod p-1$, yang \omterm{def:b3:groups:derived}{terselesaikan} jika dan hanya jika $k$ genap, karena $p - 1$
genap). Dan $x^{(p-1)/2} = g^{k(p-1)/2} = 1$ jika dan hanya jika $(p-1) \mid
k\frac{p-1}2$, yaitu jika dan hanya jika $k$ genap: kedua syaratnya bersesuaian. Untuk $x =
-1 = g^{(p-1)/2}$: ia kuadrat jika dan hanya jika $\frac{p-1}2$ genap, yaitu jika
$p \equiv 1 \pmod 4$ --- \cref{pb:b3:rings:1} sekali lagi.
\end{solution}

\begin{solution}{exo:b3:galois:5}
Di dalam $\bar{\mathbb F}_p$, $\mathbb F_{p^k} = \{x : x^{p^k} =
x\}$ merupakan himpunan tetap dari $F^k$. Irisan $\mathbb F_{p^m}
\cap \mathbb F_{p^n}$ ditetapkan oleh $F^m$ dan $F^n$, sehingga oleh
$F^{\gcd(m,n)}$ (sebab $\gcd = am + bn$: pada unsur tetap, $F^{am +
bn} = (F^m)^a(F^n)^b$ bekerja secara trivial --- eksponennya boleh diambil
positif menurut keperiodikannya); jadi irisan itu terletak di $\mathbb
F_{p^{\gcd(m,n)}}$, yang sebaliknya termuat di keduanya
(\cref{thm:b3:galois:finitefields}(3)). Untuk komposit $\mathbb
F_{p^m}\mathbb F_{p^n}$: setiap lapangan yang memuat keduanya berderajat
habis dibagi $m$ dan $n$, sehingga oleh $\operatorname{lcm}(m,n)$; dan
$\mathbb F_{p^{\operatorname{lcm}}}$ memuat keduanya: jadi itulah
kompositnya. Untuk $m \mid n$: $\operatorname{Gal}(\mathbb
F_{p^n}/\mathbb F_{p^m})$ terdiri atas pangkat $F$ yang menetapkan
$\mathbb F_{p^m}$, yakni pangkat $F^m$: jadi siklik berorde $n/m$,
dibangun oleh $F^m \colon x \mapsto x^{p^m}$ (ordenya seperti pada
\cref{thm:b3:galois:finitefields}(2)).
\end{solution}

\begin{solution}{exo:b3:galois:6}
Polinomial $X^{q^n} - X$ bersifat \omterm{def:b3:galois:separable}{separabel} (turunannya $-1$) dengan himpunan akar
$\mathbb F_{q^n}$. Misalkan $P$ monik \omterm{def:b3:rings:divisibility}{tak tereduksi} berderajat $d$.
Jika $d \mid n$: $\mathbb F_q[X]/(P) \cong \mathbb F_{q^d}
\subseteq \mathbb F_{q^n}$, sehingga $P$ berakar $\alpha \in
\mathbb F_{q^n}$; lalu $\alpha^{q^n} = \alpha$, dan $P =
\pi_\alpha$ membagi $X^{q^n} - X$. Jika $P \mid X^{q^n} - X$: sebuah
akar $\alpha \in \mathbb F_{q^n}$ membangun $\mathbb F_{q^d}
\subseteq \mathbb F_{q^n}$, sehingga $d \mid n$
(\cref{thm:b3:galois:finitefields}(3)). \omterm{def:b3:rings:divisibility}{Unsur tak tereduksi} yang berbeda
saling relatif prima dan hasil kalinya \omterm{def:b3:galois:separable}{separabel}: jadi setiap $P$ muncul dengan
eksponen tepat $1$, dan setiap akar $X^{q^n}-X$ merupakan akar
\omterm{def:b3:galois:algebraic}{polinomial minimalnya}: jadi faktorisasinya berlaku. Dengan membandingkan
\omterm{def:b3:galois:extension}{derajatnya}: $q^n = \sum_{d\mid n} d\,I_d(q)$.

Akibatnya $I_1 = q$; lalu $q^2 = I_1 + 2I_2$ memberikan $I_2 =
\frac{q^2 - q}2$; $q^3 = I_1 + 3I_3$ memberikan $I_3 = \frac{q^3 -
q}3$; dan $q^4 = I_1 + 2I_2 + 4I_4$ memberikan $I_4 = \frac{q^4 -
q^2}4$. Keberadaannya: $nI_n = q^n - \sum_{d \mid n,\, d < n}
dI_d \geq q^n - \sum_{d \leq n/2} q^d > q^n - q^{n/2 + 1} \geq
0$ untuk $n \geq 2$ (dan $I_1 = q \geq 1$): jadi $I_n \geq 1$ selalu.
\end{solution}

\begin{solution}{exo:b3:galois:7}
$L = \Q(\sqrt2, \sqrt3)$ adalah \omterm{thm:b3:galois:splitting}{lapangan pemecah} $(X^2 -
2)(X^2 - 3)$, yang \omterm{def:b3:galois:separable}{separabel}: jadi \omterm{def:b3:galois:galois}{Galois} berderajat $4$
(\cref{exo:b3:galois:1}). Sebuah automorfisma mengirim $\sqrt2 \mapsto
\pm\sqrt2$ dan $\sqrt3 \mapsto \pm\sqrt3$: jadi paling banyak $4$ pilihan,
dan $\abs G = 4$ mewujudkan semuanya: $G \cong (\Z/2\Z)^2$, dengan
unsur $\mathrm{id}, \sigma (\sqrt2 \mapsto -\sqrt2), \tau
(\sqrt3\mapsto-\sqrt3), \sigma\tau$. Subgrup berorde $2$:
$\langle\sigma\rangle, \langle\tau\rangle,
\langle\sigma\tau\rangle$, dengan lapangan tetap $\Q(\sqrt3)$,
$\Q(\sqrt2)$, $\Q(\sqrt6)$ (perhatikan bahwa $\sigma\tau$ menetapkan $\sqrt6 =
\sqrt2\sqrt3$). Kisinya: $\Q$ di bawah, ketiga lapangan
kuadrat di tengah, dan $L$ di atas --- tak ada yang lain
(\cref{thm:b3:galois:fundamental}). Untuk $(X^2-2)(X^2-3)(X^2-6)$:
\omterm{thm:b3:galois:splitting}{lapangan pemecahnya} \emph{sama}, yaitu $L$ (sebab $\sqrt6 =
\sqrt2\sqrt3$), sehingga jawabannya identik: \omterm{thm:b3:galois:fundamental}{korespondensi
Galois} merupakan invarian \emph{perluasannya}, bukan invarian
polinomial yang dipilih untuk menyajikannya.
\end{solution}

\begin{solution}{exo:b3:galois:8}
(a) Grup $G$ beraksi secara setia dan transitif (karena ketaktereduksian) pada
ketiga akarnya: $G \hookrightarrow S_3$ dengan $3 \mid \abs G$: jadi $G
\cong A_3$ atau $S_3$. Setiap $\sigma \in G$ mempermutasikan $x_i$, dan
$\sigma(\delta) = \varepsilon(\sigma)\,\delta$ (sebab $\delta$ bersifat
alternasi terhadap akarnya). Jika $\Delta$ kuadrat di $\Q$: maka
$\delta \in \Q^\times$ (perhatikan $\delta \neq 0$ karena \omterm{def:b3:galois:separable}{separabel}), sehingga
$\varepsilon(\sigma) = 1$ untuk setiap $\sigma$: jadi $G \subseteq A_3$,
dan karenanya $= A_3$. Jika bukan: $\delta \notin \Q$, sehingga ada $\sigma$ dengan
$\varepsilon(\sigma) = -1$: jadi $G = S_3$. (Pada kedua kasus $\Q(\delta)
= L^{G \cap A_3}$.)

(b) Dengan $x_1 + x_2 + x_3 = 0$: $u + v = 2x_1 - (x_2 + x_3) =
3x_1$. Lagi pula
\[
uv = \sum_i x_i^2 + (j + j^2)\sum_{i<k}x_ix_k
= \Bigl(\sum x_i\Bigr)^2 - 3\sum_{i<k}x_ix_k = -3p,
\]
dengan memakai $j + j^2 = -1$ dan $\sum_{i<k}x_ix_k = p$. Lalu
\[
u^3 + v^3 = (u+v)^3 - 3uv(u+v) = 27x_1^3 + 9p\cdot 3x_1
= 27\,(x_1^3 + px_1) = -27q .
\]
Jadi $u^3, v^3$ merupakan akar $Y^2 + 27qY - 27p^3 = 0$
(dengan hasil kali $(uv)^3 = -27p^3$): $u^3 = \frac{-27q +
\sqrt{729q^2 + 108p^3}}2$, dan $x_1 = \frac{u + v}3$ dengan $v =
-3p/u$: itulah Cardano. \omterm{def:b3:groups:derived}{Keterselesaian} $S_3$ menjadi kerangkanya:
menara $\Q \subseteq \Q(\delta) \subseteq \Q(\delta, j, u)$
menambahkan lebih dahulu sebuah akar kuadrat ($\delta$, lapangan tetap $A_3$:
langkah $S_3 \to S_3/A_3$), lalu sebuah akar pangkat tiga ($u$, sebab $u^3
\in \Q(\delta, j)$: langkah $A_3 \to \{e\}$) --- deret
\omterm{def:b3:groups:derived}{komutator} $S_3 \supset A_3 \supset \{e\}$ yang menjelma.
\end{solution}

\begin{solution}{exo:b3:galois:9}
$\eta_0 + \eta_1 = \zeta_5 + \zeta_5^2 + \zeta_5^3 + \zeta_5^4 =
-1$ (sebab jumlah semua akar satuan ke-$5$ adalah $0$). Lalu $\eta_0\eta_1 =
(\zeta + \zeta^4)(\zeta^2 + \zeta^3) = \zeta^3 + \zeta^4 +
\zeta^6 + \zeta^7 = \zeta^3 + \zeta^4 + \zeta + \zeta^2 = -1$
(indeksnya modulo $5$). Jadi $\eta_0, \eta_1$ merupakan akar $Y^2 + Y
- 1$: yakni $\frac{-1 \pm \sqrt5}2$. Karena $\eta_0 =
2\cos\frac{2\pi}5 > 0$: maka $\eta_0 = \frac{\sqrt5 - 1}2$, sehingga
$\cos\frac{2\pi}5 = \frac{\sqrt5 - 1}4$, dan $\eta_1 =
\frac{-1-\sqrt5}2$. Grup
$\operatorname{Gal}(\Q(\zeta_5)/\Q) \cong (\Z/5\Z)^\times$ bersifat
siklik berorde $4$: ia mempunyai subgrup berorde $2$ yang
\emph{tunggal} ($\{\pm 1\}$, yakni $\zeta \mapsto \zeta^{\pm1}$), sehingga
$\Q(\zeta_5)$ mempunyai satu sublapangan kuadrat saja
(\cref{thm:b3:galois:fundamental}), yang memuat $\eta_0
\notin \Q$: jadi sublapangan itu $\Q(\eta_0) = \Q(\sqrt5)$. Keterlukisannya:
$\cos\frac{2\pi}5$ terletak di menara kuadrat $\Q \subseteq
\Q(\sqrt5)$, dan $\zeta_5$ satu langkah kuadrat di atasnya:
\cref{thm:b3:galois:wantzel} melukis segi lima itu.
\end{solution}

\begin{solution}{exo:b3:galois:10}
(a) Tulis $s = S^{1/p}$, $t = T^{1/p}$ (yaitu unsur sebuah
\omterm{def:b3:galois:closure}{penutup aljabar} terpilih dengan $s^p = S$, $t^p = T$). Polinomial $X^p - S$
\omterm{def:b3:rings:divisibility}{tak tereduksi} atas pecahan $K = \mathbb F_p(S, T) = (\mathbb
F_p(T))(S)$: lewat \omterm{thm:b3:rings:criteria}{Eisenstein} pada \omterm{def:b3:rings:divisibility}{unsur prima} $S$ dari
\omterm{def:b3:rings:pidufd}{DFT} $\mathbb F_p(T)[S]$ (\cref{thm:b3:rings:criteria}). Jadi
$[K(s):K] = p$; demikian pula $X^p - T$ memenuhi \omterm{thm:b3:rings:criteria}{Eisenstein} pada $T$ atas
pecahan $K(s) = \mathbb F_p(s)(T)$ --- sebab $T$ tetap prima di
$\mathbb F_p(s)[T]$ --- sehingga $[L : K(s)] = p$ dan $[L:K] =
p^2$. Untuk $\alpha \in L$: $L = K[s, t]$, sehingga $\alpha = \sum
c_{ij}s^it^j$ (dengan $c_{ij} \in K$), dan menurut morfisma \omterm{thm:b3:galois:finitefields}{Frobenius}
$\alpha^p = \sum c_{ij}^p S^iT^j \in K$.

(b) Jika $L = K(\alpha)$, maka $[K(\alpha):K] = p^2$; padahal
$\alpha^p = a \in K$ berarti $\alpha$ menghapus $X^p - a$, sehingga
$\deg\pi_\alpha \leq p < p^2$: kontradiksi. Jadi tak ada unsur primitif:
\cref{thm:b3:galois:primitive} sungguh memerlukan
separabilitas (di sini setiap $\pi_\alpha$ membagi suatu $X^p - a =
(X - \alpha)^p$: yaitu tak \omterm{def:b3:galois:separable}{separabel} murni).
\end{solution}

\begin{solution}{exo:b3:galois:11}
(a) \omterm{thm:b3:rings:criteria}{Eisenstein} pada $2$ ($2 \mid 4, 2$; $4 \nmid 2$): jadi $P$
\omterm{def:b3:rings:divisibility}{tak tereduksi}. Jika $\alpha$ sebuah akarnya, maka $[\Q(\alpha):\Q] = 5$
membagi $[L:\Q] = \abs G$ (dengan $L$ \omterm{thm:b3:galois:splitting}{lapangan pemecahnya}): lalu Cauchy
(\cref{thm:b3:groups:cauchy}) memberikan unsur berorde $5$ di
$G \leq S_5$; dan di $S_5$, hanya siklus-$5$ yang berorde $5$ (sebab ordenya
berupa KPK panjang siklusnya).

(b) Turunan $P'(x) = 5x^4 - 4$ lenyap di $\pm(4/5)^{1/4} \approx \pm
0.946$: jadi satu maksimum lokal lalu satu minimum lokal. Nilainya: $P(-2)
= -22 < 0$, $P(0) = 2 > 0$, $P(1) = -1 < 0$, $P(2) = 26 > 0$:
tiga pergantian tanda, dan paling banyak tiga akar real (karena dua titik
kritis): jadi tepat $3$ akar real, sehingga ada sepasang akar
sekawan kompleks. Ambil \omterm{thm:b3:galois:splitting}{lapangan pemecah} $L$ di dalam $\C$:
konjugasi kompleks memetakan $L$ pada dirinya sendiri (sebab ia mempermutasikan akarnya,
yang membangun $L$) dan menetapkan $\Q$, sehingga ia mendefinisikan unsur
$G$; ia menetapkan ketiga akar realnya dan menukar dua sisanya: yaitu
sebuah transposisi.

(c) Misalkan $\tau = (a\,b)$ dan $\sigma$ sebuah siklus-$5$ di $G$. Suatu
pangkat $\sigma^k$ mengirim $a$ ke $b$ (dengan $k \ne 0 \bmod 5$), dan
$\sigma^k$ tetap siklus-$5$: setelah dinamai ulang, andaikan $\sigma =
(1\,2\,3\,4\,5)$ dan $\tau = (1\,2)$. Dengan mengonjugasikannya,
$\sigma^m\tau\sigma^{-m} = (\sigma^m(1)\ \sigma^m(2))$: jadi
transposisi bersebelahan $(1\,2), (2\,3), (3\,4), (4\,5),
(5\,1)$ semuanya terletak di $G$; sedangkan transposisi bersebelahan membangun $S_5$
(sebab setiap transposisi $(i\,j)$ merupakan hasil kali transposisi bersebelahan,
dan transposisi membangun seluruhnya). Jadi $G = S_5$, yang tak \omterm{def:b3:groups:derived}{terselesaikan}
(\cref{cor:b3:groups:snnotsolvable}), dan
\cref{thm:b3:galois:solvable} menyimpulkannya: $X^5 - 4X + 2$ tidak
\omterm{def:b3:galois:radical}{terselesaikan dengan radikal}.
\end{solution}

\begin{solution}{exo:b3:galois:12}
(a) Polinomial $X^4 - 2$ \omterm{def:b3:rings:divisibility}{tak tereduksi} (\omterm{thm:b3:rings:criteria}{Eisenstein} pada $2$):
jadi $[\Q(\alpha):\Q] = 4$; dan $\iu \notin \Q(\alpha) \subseteq \R$,
sehingga $[L : \Q(\alpha)] = 2$ dan $[L:\Q] = 8$. Perluasannya bersifat
\omterm{def:b3:galois:galois}{Galois} (\omterm{thm:b3:galois:splitting}{lapangan pemecah} \omterm{def:b3:galois:separable}{polinomial separabel}: akarnya
adalah $\iu^k\alpha$), jadi $\abs G = 8$. Sebuah automorfisma mengirim
$\alpha$ ke salah satu dari keempat akarnya dan $\iu$ ke $\pm\iu$:
jadi paling banyak $8$ pemetaan, dan semuanya terwujud. Automorfisma $\sigma$ (berorde
$4$: sebab $\sigma^2(\alpha) = -\alpha$, $\sigma^4 = e$) dan $\tau$
(berorde $2$) memenuhi
\[
\tau\sigma\tau(\alpha) = \tau\sigma(\alpha) =
\tau(\iu\alpha) = -\iu\alpha = \sigma^{-1}(\alpha),
\qquad \tau\sigma\tau(\iu) = \iu\cdot(-1)(-1) = \iu ,
\]
atau lebih cermat: $\tau\sigma\tau(\iu) = \tau\sigma(-\iu) =
\tau(-\iu) = \iu = \sigma^{-1}(\iu)$. Jadi $\tau\sigma\tau =
\sigma^{-1}$: itulah penyajian $D_4$.

(b) Kesepuluh subgrup $D_4 = \langle\sigma,
\tau\rangle$ adalah: $\{e\}$; lima subgrup berorde $2$, yakni
$\langle\sigma^2\rangle$, $\langle\tau\rangle$,
$\langle\sigma^2\tau\rangle$, $\langle\sigma\tau\rangle$,
$\langle\sigma^3\tau\rangle$; tiga subgrup berorde $4$, yakni
$\langle\sigma\rangle$, $\{e, \sigma^2, \tau,
\sigma^2\tau\}$, $\{e, \sigma^2, \sigma\tau,
\sigma^3\tau\}$; dan $D_4$ sendiri. Lapangan tetapnya (\omterm{def:b3:galois:extension}{derajat} = indeks):
$\{e\} \leftrightarrow L$; subgrup berorde $2$
$\leftrightarrow$ kelima lapangan kuartiknya
\[
\langle\tau\rangle \leftrightarrow \Q(\alpha),\quad
\langle\sigma^2\tau\rangle \leftrightarrow \Q(\iu\alpha),
\quad
\langle\sigma^2\rangle \leftrightarrow \Q(\sqrt2, \iu),\quad
\langle\sigma\tau\rangle \leftrightarrow
\Q\bigl((1+\iu)\alpha\bigr),\quad
\langle\sigma^3\tau\rangle \leftrightarrow
\Q\bigl((1-\iu)\alpha\bigr) .
\]
Pemeriksaannya: $\tau$ menetapkan $\alpha$ yang real; $\sigma^2\tau$ mengirim
$\alpha \mapsto -\alpha$ dan $\iu \mapsto -\iu$, sehingga menetapkan
$\iu\alpha$; dan karena $\sigma\tau(\alpha) = \iu\alpha$,
$\sigma\tau(\iu) = -\iu$:
\[
\sigma\tau\bigl((1+\iu)\alpha\bigr) =
(1 - \iu)\,\iu\alpha = (1 + \iu)\alpha,
\qquad
\sigma^3\tau\bigl((1-\iu)\alpha\bigr) =
(1 + \iu)(-\iu)\alpha = (1 - \iu)\alpha :
\]
jadi setiap pencerminan menetapkan pembangunnya, dan lapangan tetapnya,
yang berderajat $4 =$ indeksnya, tepat berupa lapangan yang dibangunnya
(pembangunnya akar $X^4 + 8$, yang \omterm{def:b3:rings:divisibility}{tak tereduksi}). Subgrup berorde $4$
$\leftrightarrow$ ketiga lapangan kuadratnya:
$\langle\sigma\rangle \leftrightarrow \Q(\iu)$ (sebab $\sigma$
menetapkan $\iu$); $\{e, \sigma^2, \tau, \sigma^2\tau\}
\leftrightarrow \Q(\sqrt2)$ (keempatnya menetapkan $\alpha^2$ setelah
tandanya diperiksa: $\tau(\sqrt2) = \sqrt2$, $\sigma^2(\alpha^2) =
(-\alpha)^2$); dan $\{e, \sigma^2, \sigma\tau, \sigma^3\tau\}
\leftrightarrow \Q(\iu\sqrt2)$ (sebab $\sigma\tau(\iu\alpha^2) =
(-\iu)(\iu\alpha)^2 = \iu\alpha^2$).

(c) Yang \omterm{def:b3:galois:galois}{Galois} atas $\Q$ $\leftrightarrow$ \omterm{def:b3:groups:normal}{subgrup normal}
$D_4$: yakni $\{e\}$, $\langle\sigma^2\rangle$ (pusatnya), ketiga
subgrup berorde $4$, dan $D_4$ --- sehingga lapangan antara yang
\omterm{def:b3:galois:galois}{Galois} adalah $L$, $\Q(\sqrt2, \iu)$, ketiga
lapangan kuadratnya, dan $\Q$. Kelima lapangan kuartik yang ditetapkan
oleh pencerminan tak normal tidaklah \omterm{def:b3:galois:galois}{Galois}: $\Q(\alpha)$ memuat
satu akar $X^4 - 2$ tetapi tidak memuat $\iu\alpha$ (sebab ia real) ---
konjugasi dengan $\sigma$ memindahkan $\langle\tau\rangle$ ke
$\langle\sigma^2\tau\rangle$, persis seperti ia memindahkan
$\Q(\alpha)$ ke $\Q(\iu\alpha)$: jadi ketaknormalan
subgrupnya \emph{adalah} keberadaan lapangan sekawannya.
\end{solution}

\begin{solution}{pb:b3:galois:1}
\textbf{1.} Polinomial $\Phi_{17}$ \omterm{def:b3:rings:divisibility}{tak tereduksi}
(\cref{thm:b3:galois:cyclotomicirred}, atau
\cref{ex:b3:rings:eisenstein} untuk indeks prima): jadi $[L:\Q] =
\varphi(17) = 16$ dan $G \cong (\Z/17\Z)^\times$, yang siklik berorde
$16$ (\cref{thm:b3:galois:cyclic}). Pangkat $3$ modulo
$17$:
\[
3,\ 9,\ 10,\ 13,\ 5,\ 15,\ 11,\ 16,\ 14,\ 8,\ 7,\ 4,\ 12,\ 2,\
6,\ 1
\]
--- enam belas nilai berbeda: jadi $3$ membangunnya.

\textbf{2.} $G = \langle\sigma\rangle$ siklik berorde $16$;
$H_k = \langle\sigma^{2^k}\rangle$ berorde $2^{4-k}$, dan
$[H_k : H_{k+1}] = 2$. Menurut teorema fundamental
(\cref{thm:b3:galois:fundamental}), $L_k = L^{H_k}$ memenuhi
$[L_k : \Q] = [G : H_k] = 2^k$: jadi setiap $[L_{k+1}:L_k] = 2$.

\textbf{3.} Unsur $\zeta \in L = L_4$ duduk di puncak menara perluasan
kuadrat atas $\Q$: menurut \cref{thm:b3:galois:wantzel}, $\zeta$
\omterm{def:b3:galois:constructible}{dapat dilukis}; dan segi-$17$ bertitik sudut $\zeta^k$.

\textbf{4.} Automorfisma $\sigma^2$ mengalikan eksponennya dengan $9$; sedangkan
eksponen $\eta_0$ adalah pangkat genap dari $3$,
\[
\{3^{2k} \bmod 17\} = \{1, 9, 13, 15, 16, 8, 4, 2\},
\]
yaitu himpunan yang stabil terhadap perkalian dengan $9 = 3^2$; jadi $\eta_0$
(dan demikian pula $\eta_1$) ditetapkan oleh $H_1 =
\langle\sigma^2\rangle$: berarti $\eta_0, \eta_1 \in L_1$, sebuah lapangan
kuadrat. Automorfisma $\sigma$ memetakan pangkat genap ke ganjil: ia menukar $\eta_0,
\eta_1$. Karena itu $\eta_0 + \eta_1$ dan $\eta_0\eta_1$ ditetapkan oleh
seluruh $G$: jadi rasional; sehingga $\eta_0, \eta_1$ merupakan akar sebuah
polinomial kuadrat rasional.

\textbf{5.} $\eta_0 + \eta_1 = \sum_{c=1}^{16}\zeta^c = -1$.
Hasil kalinya terjabar menjadi $64$ suku $\zeta^{a + b}$, dengan $a$ di himpunan
genap dan $b$ di himpunan ganjil. Tak ada suku yang sama dengan $\zeta^0$: sebab $b = -a$
mustahil, karena $-1 = 16 = 3^8$ merupakan pangkat \emph{genap}, sehingga
$-a$ tetap berada di himpunan genap. Dengan demikian $\eta_0\eta_1 = \sum_{c \neq 0}
n_c\zeta^c$ dengan $\sum n_c = 64$; menerapkan $\sigma$ menetapkan
$\eta_0\eta_1$ (sebab ia menukar kedua faktornya) dan mempermutasikan $\zeta^c$
secara transitif atas semua $c \neq 0$, sehingga semua $n_c$ sama: yakni $n_c
= 4$ dan $\eta_0\eta_1 = 4\sum_{c\neq0}\zeta^c = -4$.

\textbf{6.} Bilangan $\eta_{0,1}$ menyelesaikan $Y^2 + Y - 4 = 0$: yakni $\frac{-1 \pm
\sqrt{17}}2$. Secara numerik, dengan memasangkan eksponen sekawannya,
$\eta_0 = 2\bigl(\cos\tfrac{2\pi}{17} + \cos\tfrac{4\pi}{17} +
\cos\tfrac{8\pi}{17} + \cos\tfrac{16\pi}{17}\bigr) \approx 1.56
> 0$: jadi $\eta_0 = \frac{-1+\sqrt{17}}2$, $\eta_1 =
\frac{-1-\sqrt{17}}2$, dan $L_1 = \Q(\eta_0) = \Q(\sqrt{17})$.

\textbf{7.} Himpunan eksponennya: untuk $\beta_0$: $\{1, 13, 16, 4\}$ =
pangkat $3^{4k}$; untuk $\beta_2$: $\{9, 15, 8, 2\}$ = $9 \times$ himpunan
itu. Gabungannya: himpunan genap, sehingga $\beta_0 + \beta_2 = \eta_0$; demikian pula
$\beta_1 + \beta_3 = \eta_1$. Perkalian dengan $13 = 3^4$
menstabilkan himpunan eksponen setiap $\beta_i$: jadi ditetapkan oleh $H_2 =
\langle\sigma^4\rangle$; sedangkan $\sigma^2$ (yaitu $\times 9$) mengirim
$\{1,13,16,4\}$ ke $\{9, 15, 8, 2\}$: jadi ia menukar $\beta_0,
\beta_2$.

\textbf{8.} Dengan menjabarkan $\beta_0\beta_2$, keenam belas jumlah
eksponennya
\[
\{1,13,16,4\} + \{9,15,8,2\} =
\{10,16,9,3,\ 5,11,4,15,\ 8,14,7,1,\ 13,2,12,6\}
\]
mencakup $1, \dots, 16$ tepat satu kali: sehingga $\beta_0\beta_2 =
\sum_{c\ne0}\zeta^c = -1$. Menerapkan $\sigma$ (yang memetakan
$\beta_0 \mapsto \beta_1$, $\beta_2 \mapsto \beta_3$: sebab
eksponennya $\times 3$): $\beta_1\beta_3 = \sigma(\beta_0\beta_2)
= -1$.

\textbf{9.} Bilangan $\beta_0, \beta_2$ menyelesaikan $Y^2 - \eta_0 Y - 1 = 0$,
sehingga $\beta_0 = \frac{\eta_0 + \sqrt{\eta_0^2 + 4}}2$ (secara numerik
$\beta_0 = 2\cos\frac{2\pi}{17} + 2\cos\frac{8\pi}{17} \approx
2.05 > 0$, jadi tanda $+$). Demikian pula $\beta_1 = \frac{\eta_1 +
\sqrt{\eta_1^2 + 4}}2 \approx 0.344$ (pemeriksaan numerik kembali menetapkan
tandanya). Lalu $L_2 = L_1(\beta_0)$, yang kuadrat atas $L_1$.

\textbf{10.} $\gamma_0 + \gamma_1 = \zeta + \zeta^{16} +
\zeta^{13} + \zeta^4 = \beta_0$. Dan
\[
\gamma_0\gamma_1 = (\zeta + \zeta^{16})(\zeta^{13} + \zeta^4)
= \zeta^{14} + \zeta^{5} + \zeta^{12} + \zeta^{3} = \beta_1 .
\]
Jadi $\gamma_0, \gamma_1$ menyelesaikan $Y^2 - \beta_0Y + \beta_1 = 0$;
secara numerik $\gamma_0 = 2\cos\frac{2\pi}{17} \approx 1.865 >
\gamma_1 \approx 0.185$: sehingga $\gamma_0 = \frac{\beta_0 +
\sqrt{\beta_0^2 - 4\beta_1}}2$.

\textbf{11.} Dengan merantainya:
\[
\eta_0 = \frac{-1 + \sqrt{17}}2,
\quad
\beta_0 = \frac{\eta_0 + \sqrt{\eta_0^2 + 4}}2,
\quad
\beta_1 = \frac{\eta_1 + \sqrt{\eta_1^2 + 4}}2,
\quad
\cos\frac{2\pi}{17} = \frac{\beta_0 + \sqrt{\beta_0^2 -
4\beta_1}}4 .
\]
Secara numerik: $\sqrt{17} \approx 4.1231$, $\eta_0 \approx
1.5616$, $\eta_1 \approx -2.5616$, $\beta_0 \approx 2.0494$,
$\beta_1 \approx 0.3441$, $\beta_0^2 - 4\beta_1 \approx
2.8234$, dan $\cos\frac{2\pi}{17} \approx \frac{2.0494 +
1.6803}4 \approx 0.93242$ --- berbanding $\cos\frac{2\pi}{17} =
0.93247\dots$: selisih kecilnya berasal dari pembulatan pada
tampilan antaranya; dengan lebih banyak angka diperoleh kembali
$0.932472$.

\textbf{12.} Konstruksinya memerlukan $[L:\Q] = p - 1$ berupa
pangkat $2$, agar rantai penuh subgrup berindeks $2$
tersedia. Jika $p = 2^m + 1$ prima dan $m = ab$ dengan $a$ ganjil $>
1$: maka $x + 1 \mid x^a + 1$ pada $x = 2^b$ menunjukkan bahwa $2^b + 1$ membagi
$p$ secara sejati --- mustahil. Jadi $m$ berupa pangkat $2$: $p =
2^{2^t} + 1$, sebuah \emph{prima Fermat} ($3, 5, 17, 257, 65537$,
\dots). Sebaliknya untuk $p$ semacam itu, $\varphi(p) = 2^{2^t}$ dan
argumen pertanyaan 1--3 (atau
\cref{cor:b3:galois:impossible}) berlaku: jadi segi-$p$ beraturan
\omterm{def:b3:galois:constructible}{dapat dilukis} jika dan hanya jika $p$ merupakan prima Fermat.

\textbf{13.} Nilai $\varphi(n)$ berupa pangkat $2$ tepat ketika $n =
2^a p_1\cdots p_r$ dengan $p_i$ berupa prima Fermat yang berbeda
(lewat kemultiplikatifan $\varphi$; sedangkan pangkat prima ganjil $p^k$ dengan $k
\geq 2$ menyumbang faktor $p \nmid 2^m$). Untuk $n \leq
20$, segi-$n$ beraturan yang \omterm{def:b3:galois:constructible}{dapat dilukis} adalah
\[
n = 3, 4, 5, 6, 8, 10, 12, 15, 16, 17, 20
\]
dengan nilai berturut-turut
\[
\varphi(n) = 2,\ 2,\ 4,\ 2,\ 4,\ 4,\ 4,\ 8,\ 8,\ 16,\ 8 .
\]
Yang mustahil adalah $n = 7, 9, 11, 13, 14, 18, 19$, dengan
$\varphi(n) = 6, 6, 10, 12, 6, 6, 18$ yang berfaktor prima ganjil.

\textbf{14.} Kuadratnya membentuk peta morfisma pengkuadratan
pada $(\Z/p\Z)^\times$ yang siklik, dengan indeks $2$:
jadi ada $\frac{p-1}2$ kuadrat dan $\frac{p-1}2$ bukan-kuadrat, sehingga
jumlah simbolnya $0$. Lalu
\[
\sum_{a=0}^{p-1}\zeta^{a^2} = 1 + 2\sum_{b \text{ kuadrat}
\neq 0}\zeta^b
= 1 + \sum_{b\neq0}\Bigl(1 + \Bigl(\frac bp\Bigr)\Bigr)
\zeta^b = \sum_{b}\zeta^b + g = g,
\]
dengan memakai $1 + \bigl(\frac bp\bigr) = \#\{a : a^2 = b\}$ dan
$\sum_{b=0}^{p-1}\zeta^b = 0$. Untuk $p = 17$: kuadrat modulo
$17$ adalah pangkat genap dari pembangun $3$, yakni
eksponen yang muncul di $\eta_0$ (Bagian II), sehingga $g =
\sum_{\text{genap }k}\zeta^{3^k} - \sum_{\text{ganjil }k}
\zeta^{3^k} = \eta_0 - \eta_1$.

\textbf{15.} Dengan $b = c - a$ ($a, b$ menjelajahi sisa tak nol,
sedangkan $c = a + b$ menjelajahi semua sisa):
\[
g^2 = \sum_c\zeta^c\sum_{a\neq0,\,a\neq c}
\Bigl(\frac{a(c-a)}p\Bigr) .
\]
Untuk $c = 0$: $\bigl(\frac{-a^2}p\bigr) =
\bigl(\frac{-1}p\bigr)$, yang dijumlahkan atas $p - 1$ nilai. Untuk $c \neq 0$: substitusikan $c - a = at$, yakni $t = c/a -
1$; ketika $a$ menjelajahi sisa tak nol, $t$ menjelajahi
secara bijektif sisa $\neq -1$ (baliklah: $a = c/(1 +
t)$). Sukunya menjadi $\bigl(\frac{a^2t}p\bigr) =
\bigl(\frac tp\bigr)$, dan
\[
\sum_{t \neq -1}\Bigl(\frac tp\Bigr)
= -\Bigl(\frac{-1}p\Bigr)
\]
(sebab jumlah penuhnya lenyap menurut pertanyaan 14). Karena itu
\[
g^2 = \Bigl(\frac{-1}p\Bigr)\Bigl[(p-1) -
\sum_{c\neq0}\zeta^c\Bigr]
= \Bigl(\frac{-1}p\Bigr)\,p = p^* ,
\]
dengan memakai $\sum_{c\neq0}\zeta^c = -1$ dan kriteria Euler
$\bigl(\frac{-1}p\bigr) = (-1)^{(p-1)/2}$. Untuk $p = 17$:
$(\eta_0 - \eta_1)^2 = (\eta_0 + \eta_1)^2 - 4\eta_0\eta_1 =
1 + 16 = 17$, cocok dengan Bagian II.

\textbf{16.} Kesamaan $g^2 = p^*$ menampilkan $\sqrt{p^*} = \pm g \in
\Q(\zeta_p)$, sehingga $\Q(\sqrt{p^*})$ menjadi sublapangan kuadrat.
Ketunggalannya: sublapangan berderajat $2$ bersesuaian, lewat \omterm{thm:b3:galois:fundamental}{korespondensi
Galois}, dengan subgrup berindeks $2$ dari
$\operatorname{Gal}(\Q(\zeta_p)/\Q) \cong (\Z/p\Z)^\times$ yang siklik,
dan grup siklik berorde genap mempunyai tepat satu subgrup
semacam itu (yaitu kuadratnya). Setiap lapangan kuadrat berbentuk
$\Q(\sqrt{d})$ dengan $d$ bebas kuadrat, dan menggabungkan lapangan
$\Q(\sqrt{p^*})$, $\Q(\iu) \subseteq \Q(\zeta_4)$ dan
$\Q(\sqrt2) \subseteq \Q(\zeta_8)$ di dalam satu
$\Q(\zeta_N)$ menangkap setiap $\sqrt d$: itulah kasus kuadrat dari
Kronecker--Weber.

\textbf{17.} Di ring komutatif mana pun, $(x + y)^q = x^q + y^q
+ q(\cdots)$: sebab koefisien binomial $\binom qk$ dengan $0 < k <
q$ habis dibagi bilangan prima $q$. Dengan mengiterasikannya pada $p - 1$
suku $g$:
\[
g^q \equiv \sum_{a}\Bigl(\frac ap\Bigr)^{q}\zeta^{aq}
= \sum_a\Bigl(\frac ap\Bigr)\zeta^{aq}
\pmod{q\Z[\zeta]},
\]
($q$ ganjil: simbolnya tak berubah). Indeks ulang $b = aq$: maka $a =
q^{-1}b$ dan $\bigl(\frac{q^{-1}b}p\bigr) =
\bigl(\frac{q}p\bigr)\bigl(\frac bp\bigr)$ (lewat kemultiplikatifan;
sebab $\bigl(\frac{q^{-1}}p\bigr) = \bigl(\frac qp\bigr)$ karena
simbol sebuah invers sama dengan simbolnya): jadi $g^q \equiv
\bigl(\frac qp\bigr)g$.

\textbf{18.} Kesamaan $g^q = g\,(g^2)^{(q-1)/2} = g\,(p^*)^{(q-1)/2}$
berlaku persis (pertanyaan 15), dan kriteria Euler di $\Z$ memberikan
$(p^*)^{(q-1)/2} \equiv \bigl(\frac{p^*}q\bigr) \pmod q$,
sehingga juga modulo $q\Z[\zeta]$: $g^q \equiv \bigl(\frac{p^*}q\bigr)g$.
Dengan membandingkannya dengan pertanyaan 17 lalu mengalikannya dengan $g$:
\[
\Bigl(\frac qp\Bigr)p^* \equiv \Bigl(\frac{p^*}q\Bigr)p^*
\pmod{q\Z[\zeta]} .
\]
Kedua ruasnya bilangan bulat rasional; selisihnya, yang $0$ atau
$\pm2p^*$, terletak di $q\Z[\zeta] \cap \Z = q\Z$ (sebab bilangan bulat $m
\in q\Z[\zeta]$ mempunyai $m/q \in \Q \cap \Z[\zeta] = \Z$, dan
yang terakhir itu karena $1, \zeta, \dots, \zeta^{p-2}$ merupakan
basis-$\Q$ yang koordinat rasionalnya menunjukkan
kebulatannya). Karena $q \nmid 2p^*$ ($q$ ganjil dan $q \neq p$),
selisihnya $0$: jadi $\bigl(\frac qp\bigr) =
\bigl(\frac{p^*}q\bigr)$.

\textbf{19.} Lewat kemultiplikatifan, $\bigl(\frac{p^*}q\bigr) =
\bigl(\frac{-1}q\bigr)^{(p-1)/2}\bigl(\frac pq\bigr) =
(-1)^{\frac{q-1}2\cdot\frac{p-1}2}\bigl(\frac pq\bigr)$,
sehingga pertanyaan 18 terbaca $\bigl(\frac qp\bigr)\bigl(\frac
pq\bigr) = (-1)^{\frac{p-1}2\frac{q-1}2}$: itulah resiprositasnya.
Periksa $(17, 3)$: eksponennya $\frac{16}2\cdot\frac22 = 8$
bernilai genap, sehingga kedua simbolnya harus bersesuaian; kuadrat modulo $3$ adalah $\{1\}$ sedangkan $17 \equiv 2$:
jadi $\bigl(\frac{17}3\bigr) = -1$; sementara kuadrat modulo $17$ adalah $\{1, 4,
9, 16, 8, 2, 15, 13\}$ dan $3$ tak ada di sana:
jadi $\bigl(\frac3{17}\bigr) = -1$. Hasil kalinya $+1$, seperti diramalkan. Untuk $x^2
\equiv 219 \pmod{383}$: $\bigl(\frac{219}{383}\bigr) =
\bigl(\frac3{383}\bigr)\bigl(\frac{73}{383}\bigr)$. Pertama:
$383 \equiv 3 \pmod4$ dan $3 \equiv 3$: resiprositas memberikan
$\bigl(\frac3{383}\bigr) = -\bigl(\frac{383}3\bigr) =
-\bigl(\frac23\bigr) = -(-1) = +1$. Kedua: $73 \equiv 1
\pmod 4$: sehingga $\bigl(\frac{73}{383}\bigr) =
\bigl(\frac{383}{73}\bigr) = \bigl(\frac{18}{73}\bigr) =
\bigl(\frac2{73}\bigr)$ (sebab $18 = 2\cdot3^2$), dan $73 \equiv 1
\pmod 8$ membuat $2$ menjadi kuadrat modulo $73$ (lewat hukum pelengkapnya,
yang dapat dibuktikan dengan $g = \zeta_8 + \zeta_8^{-1} = \sqrt2$ di
$\Q(\zeta_8)$ lewat metode yang sama): jadi $+1$. Totalnya $+1$: sehingga
kekongruenannya \omterm{def:b3:groups:derived}{terselesaikan}.

\textbf{20.} Keberadaan dan ketunggalan akar primitifnya:
jika $w$ berhimpunan periode $\{d : w = u^{n/d},\ \abs u = d\}$,
maka $d_0$ terkecil di antaranya membagi setiap periode $d$ lainnya
(sebab jika $w$ sekaligus pangkat-$d$ dan pangkat-$d'$, maka ia
pangkat-$\gcd(d, d')$: bandingkan hurufnya pada indeks yang bersesuaian
modulo FPB-nya, lewat Bézout), dan blok berpanjang $d_0$ itu
primitif. Dengan memilah $q^n$ kata menurut panjang akar
primitifnya: $q^n = \sum_{d \mid n}A(d)$. Karena $A$ dan
$d\,N_q(d)$ memenuhi rekursi yang sama dengan nilai yang sama
pada $n = 1$ (keduanya saling menentukan secara induktif dari $q^n
= \sum_{d\mid n}(\cdot)$), keduanya sama: jadi $A(d) =
d\,N_q(d)$. Bijeksinya: unsur $\alpha$ berderajat $d$
melahirkan kata $w_\alpha$ berisi koefisien\dots\
lebih baik secara langsung: unsur berderajat $d$ di
$\overline{\mathbb F_q}$ adalah akar dari $N_q(d)$
polinomial \omterm{def:b3:rings:divisibility}{tak tereduksi} berderajat $d$, yang masing-masing menyumbang $d$
akar berbeda (karena \omterm{def:b3:galois:separable}{separabel}): jadi ada $d\,N_q(d)$ unsur
berderajat $d$, cocok dengan cacahan $q^n = \sum_{d\mid n}\#\{
\text{unsur berderajat } d \text{ di } \mathbb F_{q^n}\}$
--- ayakan yang sama, sekali pada kata, sekali pada unsur lapangan.

\textbf{21.} Lemanya: $\sum_{d \mid m}\mu(d) = \mathbf
1_{m=1}$. Untuk $m > 1$ tetapkan sebuah bilangan prima $p \mid m$: pembagi
$m$ yang bebas kuadrat berpasangan sebagai $\{d, pd\}$ dengan $p \nmid d$,
dan $\mu(pd) = -\mu(d)$: sehingga jumlahnya saling meniadakan. Lalu, untuk $f(n) =
\sum_{e\mid n}g(e)$:
\[
\sum_{d \mid n}\mu(d)\,f\bigl(\tfrac nd\bigr)
= \sum_{d \mid n}\mu(d)\!\!\sum_{e \mid n/d}\!\!g(e)
= \sum_{e \mid n}g(e)\!\!\sum_{d \mid n/e}\!\!\mu(d)
= g(n) .
\]
Dengan $f(n) = q^n$ dan $g(n) = nN_q(n)$
(\cref{exo:b3:galois:6}): diperoleh $N_q(n) =
\frac1n\sum_{d\mid n}\mu(d)\,q^{n/d}$.

\textbf{22.} Suku $d = 1$ bernilai $q^n$; sedangkan setiap suku lain mempunyai
$\abs{\mu(d)q^{n/d}} \leq q^{n/2}$, dan secara kasar
$\sum_{d \mid n, d > 1}q^{n/d} \leq \sum_{j \leq
n/2}q^j < 2q^{n/2}$ (deret geometri, dengan $q \geq 2$). Jadi $nN_q(n) >
q^n - 2q^{n/2} \geq 0$ untuk $n \geq 1$: sehingga \omterm{def:b3:rings:divisibility}{unsur tak tereduksi} dengan
\omterm{def:b3:galois:extension}{derajat} apa pun memang ada, dan $\mathbb F_q[X]/(\pi) = \mathbb F_{q^n}$
dibangun kembali --- keberadaan yang disertai sensus. Proporsi
polinomial \omterm{def:b3:rings:divisibility}{tak tereduksi} di antara $q^n$ polinomial monik berderajat $n$ adalah
$\frac1n(1 + O(q^{-n/2}))$: itulah teorema bilangan prima untuk
$\mathbb F_q[X]$, dengan $n$ berperan sebagai $\log x$. Untuk $q = 2$,
cacahan $2, 1, 2, 3$ pada \cref{exo:b3:galois:6} cocok dengan
rumus itu: misalnya $N_2(4) = \frac14(2^4 - 2^2) = 3$.

\textbf{23.} Di dalam grup abelian multiplikatif berisi fungsi
rasional tak nol atas $\mathbb F_q$, tetapkan $F(n) = X^{q^n} -
X$ dan $G(n) = \prod_{\deg\pi = n}\pi$; lalu
\cref{exo:b3:galois:6} mengatakan $F(n) = \prod_{d\mid n}G(d)$.
Argumen Möbius pada pertanyaan 21, yang ditulis
secara multiplikatif (di sini eksponennya menjumlah persis seperti jumlah tadi),
memberikan $G(n) = \prod_{d \mid n}F(d)^{\mu(n/d)}$. Untuk $q = 2$,
$n = 2$: $G(2) = \frac{X^4 - X}{X^2 - X} = \frac{X(X^3 -
1)}{X(X - 1)} = X^2 + X + 1$, yakni satu-satunya polinomial kuadrat
\omterm{def:b3:rings:divisibility}{tak tereduksi} atas $\mathbb F_2$, sebagaimana seharusnya.

\textbf{24.} Karena $\omega^2 = \iu$ dan $\omega^{-2} = -\iu$, maka
$g^2 = \omega^2 + 2 + \omega^{-2} = 2$. Mimpi mahasiswa baru di
ring komutatif $\Z[\omega]/q\Z[\omega]$: $g^q =
(\omega + \omega^{-1})^q \equiv \omega^q + \omega^{-q}$. Nilai
$\omega^q + \omega^{-q}$ hanya bergantung pada $q \bmod 8$:
untuk $q \equiv \pm1$, $\omega^q + \omega^{-q} = \omega^{\pm1} +
\omega^{\mp1} = g$; sedangkan untuk $q \equiv \pm3$, dengan memakai $\omega^4 =
-1$, $\omega^{3} = -\omega^{-1}$ dan $\omega^{-3} = -\omega$,
sehingga $\omega^q + \omega^{-q} = -g$. Di pihak lain $g^q =
g\,(g^2)^{(q-1)/2} = g\,2^{(q-1)/2} \equiv \bigl(\frac2q\bigr)
g \pmod{q\Z[\omega]}$ menurut kriteria Euler modulo $q$. Dengan membandingkan
lalu mengalikannya dengan $g$: $2\bigl(\frac2q\bigr) \equiv
\pm2 \pmod{q\Z[\omega]}$; seandainya tandanya berbeda, $q$ akan
membagi $4$ di $\Z[\omega]$, sehingga juga di $\Z$ (sebab $q\Z[\omega] \cap
\Z = q\Z$: lihat koordinatnya pada basis $1, \omega, \omega^2,
\omega^3$), dan itu mustahil untuk $q$ ganjil. Jadi $\bigl(\frac2q\bigr) =
+1$ jika dan hanya jika $q \equiv \pm1 \pmod 8$. Pemeriksaan keparitasan: $q = 8k \pm 1$
memberikan $(q^2 - 1)/8 = 2k(4k \pm 1)$, yang genap; sedangkan $q = 8k \pm 3$
memberikan $(q^2 - 1)/8 = 8k^2 \pm 6k + 1$, yang ganjil: jadi rumus
$(-1)^{(q^2-1)/8}$ mengemas kedua kasus itu. Secara numerik: $3^2
= 9 \equiv 2 \pmod 7$ (sebab $7 \equiv -1$), $6^2 = 36 \equiv 2
\pmod{17}$ (sebab $17 \equiv 1$); sedangkan kuadrat modulo $3$ adalah $\{0,
1\}$ dan modulo $5$ adalah $\{0, 1, 4\}$, dan tak satu pun memuat $2$
(sebab $3 \equiv 3$, $5 \equiv -3 \pmod 8$).

\textbf{25.} Setiap $f \in \mathbb F_q[X]$ yang monik terfaktorkan
secara tunggal sebagai $\prod_\pi\pi^{e_\pi}$ atas polinomial monik
\omterm{def:b3:rings:divisibility}{tak tereduksi}: dengan memilahnya menurut \omterm{def:b3:galois:extension}{derajat},
\[
\sum_{f \text{ monic}}t^{\deg f}
= \prod_\pi\ \sum_{e \geq 0}t^{e\deg\pi}
= \prod_\pi\bigl(1 - t^{\deg\pi}\bigr)^{-1}
= \prod_{n \geq 1}\bigl(1 - t^n\bigr)^{-N_q(n)},
\]
dan semua hasil kali itu sah secara t-adik (sebab hanya \omterm{def:b3:galois:extension}{derajat} $\leq m$
yang menyentuh koefisien $t^m$, dan polinomial \omterm{def:b3:rings:divisibility}{tak tereduksi} pada setiap \omterm{def:b3:galois:extension}{derajat}
berhingga banyaknya). Ruas kirinya sama dengan $\sum_m
q^mt^m = (1 - qt)^{-1}$: itulah identitasnya. Logaritmanya: $-\log(1 -
qt) = \sum_m\frac{q^m}mt^m$, sedangkan $\sum_nN_q(n)\bigl(-\log(1
- t^n)\bigr) = \sum_nN_q(n)\sum_k\frac{t^{nk}}k$; koefisien
$t^m$ memberikan $\frac{q^m}m = \sum_{dk = m}
\frac{N_q(d)}k = \frac1m\sum_{d \mid m}d\,N_q(d)$, yakni
$q^m = \sum_{d\mid m}d\,N_q(d)$. Pemeriksaan dengan tangan, $q = 2$,
untuk koefisien $t^2$: $N_2(1) = 2$, $N_2(2) = 1$, dan
$(1 - t)^{-2}(1 - t^2)^{-1} = (1 + 2t + 3t^2 + \dots)(1 + t^2
+ \dots)$ berkoefisien $t^2$ sebesar $3 + 1 = 4 = 2^2$. Akhirnya,
rumus pertanyaan 21 dengan pembagi $1, 2, 3, 6$:
\[
N_2(6) = \tfrac16\bigl(2^6 - 2^3 - 2^2 + 2\bigr)
= \tfrac{54}6 = 9,
\]
dan memang $1\cdot2 + 2\cdot1 + 3\cdot2 + 6\cdot9 = 2 + 2 +
6 + 54 = 64 = 2^6$.

\end{solution}
